% sec_B.tex -- A bulk estimate uniform in the vision radius (Theorems R and B, conditional Corollary M).
% Source: referee/P6a-claims.md, version 4 (stream P6a, rounds 6--6d), which reconciles the fixes of
% reports/referee-RP6a-1.md and reports/referee-RP6a-2.md. Doubly refereed, except Remark B-relaxed (added later).
% Corollary M is CONDITIONAL: (ST) is proved in sec_ST.tex; (T) is not proved.

\newcommand{\Bjs}{j_{\mathrm s}}
\newcommand{\Bts}{t_{\mathrm s}}
\newcommand{\BLs}{L_{\mathrm s}}
\newcommand{\BTs}{T_{\mathrm s}}
\newcommand{\BF}{\mathcal F}

\section{A bulk estimate uniform in the vision radius}\label{sec:bulk}

Theorem~\ref{thm:D2} loses a logarithm because it bounds the stretching rate of the virtual front by $v/\rho$ instead of
$v\sin^2\!\phi/\rho$ (Remark~\ref{rem:D2}). This section controls the misalignment for a greedy chase of the virtual front with
step $\varepsilon$, away from the two ends of the sweep. Theorem~\ref{thm:R} reduces an $\varepsilon$-independent sweep time to three
properties of the front seen from the current base point: flatness of the bulk \textup{(F$_b$)}, separation \textup{(S)} and a bound
on the tip \textup{(T)}. Theorem~\ref{thm:Bulk} proves that \textup{(F$_b$)} and \textup{(S)} propagate along the whole sweep from a
start condition \textup{(ST)} at one leg, with constants that do not depend on $\varepsilon$, for $v\le10^{-12}$.
Corollary~\ref{cor:M} combines the two. \textbf{It is conditional}: \textup{(ST)} is proved for an explicit start in
Section~\ref{sec:start} (Theorem~\ref{thm:ST}), but \textup{(T)} is not proved (Remark~\ref{rem:B-status}).
An unconditional $\varepsilon$-independent sweep is obtained in Section~\ref{sec:main} by a different route (a tangential start and
Theorem~\ref{thm:End} in place of Theorem~\ref{thm:R}); it uses the lemmas of this section but neither Corollary~\ref{cor:M} nor
\textup{(T)}. The constants are deliberately wasteful.

Normalise $d=1$, $0<\varepsilon<1$, $0<v<1$, and put $\mu:=v/(1-v)$. Identify $\R^2$ with $\mathbb C$,
$a\cdot b:=\operatorname{Re}(\bar ab)$, $a\times b:=\operatorname{Im}(\bar ab)=a_1b_2-a_2b_1$, and $a^\perp:=ia$ (rotation by
$+90^\circ$), so that $a^\perp\times b=-a\cdot b$. External inputs: Proposition~\ref{prop:length}(2), Lemma~\ref{lem:A41} (only in
Corollary~\ref{cor:M}), and $C^1$ dependence on initial data for an ordinary differential equation whose field is $C^1$ with bounded
derivative.

\subsection{The field and the virtual candidates}

We use a $C^1$ variant of the field $x/\max(|x|,\varepsilon)$ of Lemma~\ref{lem:D2-field}. Let
\[
f(\rho):=\begin{cases}1,&\rho\ge\varepsilon,\\ (\rho/\varepsilon)(2-\rho/\varepsilon),&0\le\rho\le\varepsilon,\end{cases}\qquad
F(x):=f(|x|)\frac{x}{|x|}\ \ (F(0):=0),\qquad A(x):=DF(x),
\]
and $a^*(\rho):=1/\rho$ for $\rho\ge\varepsilon$, $a^*(\rho):=2/\varepsilon$ for $\rho<\varepsilon$.

\begin{lemma}[The field]\label{lem:B-V}
\begin{enumerate}
\item[(i)] $F\in C^1(\R^2)$, $|F|\le1$, and $F(x)=x/|x|$ for $|x|\ge\varepsilon$.
\item[(ii)] With $\rho=|x|$, $\rhat=x/\rho$, $\Pi:=I-\rhat\rhat^{\mathsf T}$:
$A=f'(\rho)\,\rhat\rhat^{\mathsf T}+(f(\rho)/\rho)\,\Pi$; $A$ is symmetric and $0\le A\le a^*(\rho)I$; for $\rho\ge\varepsilon$,
$A=\Pi/\rho$.
\item[(iii)] For every unit vector $e$: $|e^\perp\cdot A(x)e|\le|x\times e|/\max(\rho,\varepsilon)^2$.
\item[(iv)] $(F(x)-F(y))\cdot(x-y)\ge0$.
\end{enumerate}
\end{lemma}
\begin{proof}
$F(x)=g(|x|)x$ with $g(\rho)=(2\varepsilon-\rho)/\varepsilon^2$ on $[0,\varepsilon]$ and $g=1/\rho$ on $[\varepsilon,\infty)$. $g$ is
continuous, and $g'=-1/\varepsilon^2$, resp.\ $-1/\rho^2$, agree at $\rho=\varepsilon$, so $g\in C^1(0,\infty)$; near $0$,
$F(x)=(2\varepsilon x-|x|x)/\varepsilon^2$ and $x\mapsto|x|x$ is $C^1$. This gives (i) ($|F|=f\le1$). Next
$DF=gI+\rho g'\rhat\rhat^{\mathsf T}=(f/\rho)\Pi+(g+\rho g')\rhat\rhat^{\mathsf T}$ and $g+\rho g'=f'$. Inside the disc
$f'=(2/\varepsilon)(1-\rho/\varepsilon)\in[0,2/\varepsilon]$ and $f/\rho=(2-\rho/\varepsilon)/\varepsilon\in[1/\varepsilon,2/\varepsilon]$;
outside $f'=0$, $f/\rho=1/\rho$. This is (ii). (iii) Write $e=\cos\chi\,\rhat+\sin\chi\,n$ with $n:=\rhat^\perp$; then
$e^\perp=-\sin\chi\,\rhat+\cos\chi\,n$ and $e^\perp\cdot Ae=\sin\chi\cos\chi\,(f/\rho-f')$, where $f/\rho-f'=\rho/\varepsilon^2$ for
$\rho<\varepsilon$ and $=1/\rho$ for $\rho\ge\varepsilon$, i.e.\ $=\rho/\max(\rho,\varepsilon)^2$; and $|\sin\chi|=|x\times e|/\rho$.
(iv) $(F(x)-F(y))\cdot(x-y)=\int_0^1(x-y)\cdot A(y+\sigma(x-y))(x-y)\,d\sigma\ge0$.
\end{proof}
$F$ has a Lipschitz derivative but is not $C^2$ ($A$ is only piecewise smooth across $\rho=\varepsilon$); nothing below uses more than
$C^1$, except Lemma~\ref{lem:B-1}(iii),(iv), which is stated where the front is $C^2$.

Let $P$ be a $1$-Lipschitz path. The \emph{virtual candidates} are defined as in Lemma~\ref{lem:D2-virtual}, with the present field:
$\tilde S_\theta'=vF(\tilde S_\theta-P(t))$, $\tilde S_\theta(0)=e^{i\theta}$. In this section we abbreviate
\[
x_\theta:=\tilde S_\theta,\qquad r_\theta:=x_\theta-P,\qquad\rho_\theta:=|r_\theta| .
\]

\begin{lemma}[Virtual candidates]\label{lem:B-VC}
\begin{enumerate}
\item[(a)] $x_\theta$ exists and is unique on $[0,\infty)$, $|x_\theta'|\le v$; $\theta\mapsto x_\theta(t)$ is $C^1$,
$\partial_t\partial_\theta x=vA(r)\partial_\theta x$, and $|\partial_\theta x(\theta,t)|$ is non-decreasing in $t$ and $\ge1$.
\item[(b)] $|x_\theta(t)-x_\phi(t)|$ is non-decreasing in $t$.
\item[(c)] If $\rho_\theta(t)\le\varepsilon$ for some $t$, the true candidate $\theta$ is seen by time $t$.
\end{enumerate}
\end{lemma}
\begin{proof}
(a) $vF(\cdot-P(t))$ is continuous in $t$ and $C^1$ in $x$ with $\|DF\|\le2/\varepsilon$: global existence, uniqueness, $C^1$
dependence on the initial point, and the variational equation. Then
$\partial_t|\partial_\theta x|^2=2v\,\partial_\theta x\cdot A\,\partial_\theta x\ge0$ and $|\partial_\theta e^{i\theta}|=1$.
(b) $\frac{d}{dt}|x_\theta-x_\phi|^2=2v(F(r_\theta)-F(r_\phi))\cdot(r_\theta-r_\phi)\ge0$. (c) As in
Lemma~\ref{lem:D2-virtual}(c): let $t_1:=\inf\{s:\rho_\theta(s)\le\varepsilon\}\le t$ (note $\rho_\theta(0)=1>\varepsilon$). On
$[0,t_1]$, $\rho_\theta\ge\varepsilon$, where $F$ is the unit field, so $x_\theta$ solves the true equation there; by uniqueness it is
the true candidate on $[0,t_1]$, and its distance to Harry at $t_1$ is $\le\varepsilon$.
\end{proof}

\emph{Notation for the front.} Fix $t$. The front $\theta\mapsto x_\theta(t)$ is a regular $C^1$ curve; $ds=|\partial_\theta x|d\theta$
is its arclength and $T$ its unit tangent. For labels $a<b$ put
\[
l(a,b):=\int_a^bds,\qquad c(a,b):=|x_b-x_a|,\qquad e=e_{ab}:=\frac{x_b-x_a}{c},
\]
\[
E(a,b):=l-c\ge0,\qquad N(a,b):=(2E)^{1/2}.
\]
Here $c(a,b)>0$ for $b-a<2\pi$ by Lemma~\ref{lem:B-VC}(b) (while $c(0,2\pi)=0$). $E$ is the \emph{two-point excess}. A subscript
$t$ or a third argument, as in $l_t(a,b)$, $c(a,b;t)$, indicates the time. At a point of the front, $\rho$ and $\rhat$ refer to
$r=x-P$, and $\phi$ is the angle from $T$ to $\rhat$: $\cos\phi:=T\cdot\rhat$, $\sin\phi:=T\times\rhat$.

\subsection{Evolution laws}

\begin{lemma}[Pointwise laws]\label{lem:B-1}
Let $w:=|\partial_\theta x|$.
\begin{enumerate}
\item[(i)] $w_t=vw\,T\cdot AT$; for $\rho>\varepsilon$ this is $w_t/w=v\sin^2\!\phi/\rho$.
\item[(ii)] $T_t=v(AT-(T\cdot AT)T)$; for $\rho>\varepsilon$ the tangent angle $\psi$ ($T=e^{i\psi}$) obeys
$\psi_t=-(v/\rho)\sin\phi\cos\phi$.
\item[(iii)] Let $\theta_0,t_0$ be such that $\rho>\varepsilon$ on a neighbourhood and $x(\cdot,t)$ is $C^2$ there (true whenever the
points of a neighbourhood of $\theta_0$ have stayed in $\{\rho>\varepsilon\}$ since time $0$). With the curvature $k:=\partial_s\psi$:
$\partial_s\rho=\cos\phi$, $\partial_s\phi=-k-\sin\phi/\rho$, and
\[
k_t=\frac v\rho(1-3\sin^2\!\phi)\,k+\frac v{\rho^2}\sin\phi\,(3\cos^2\!\phi-1).
\]
\item[(iv)] If moreover $P(t)$ lies on the front, at arclength $0$, and the arc $[0,s]$ is $C^2$ with $\rho>0$ on $(0,s]$, then
$\rho\sin\phi\,(s)=-\int_0^sk\rho\cos\phi\,du$, hence
\[
k_t(s)=\frac v\rho\Bigl[(1-3\sin^2\!\phi)\,k(s)-(3\cos^2\!\phi-1)\rho^{-2}\int_0^sk\rho\cos\phi\,du\Bigr].
\]
\end{enumerate}
\end{lemma}
So, to leading order in $\phi$, $k_t=(v/\rho)(k(s)-\text{a weighted mean of $k$ on }[0,s])$: constant curvature is neutral. Parts
(ii)--(iv) are not used in this section; (ii) and (iii) are used in later sections.
\begin{proof}
(i), (ii): $\partial_t\partial_\theta x=vA\partial_\theta x$ and $\partial_\theta x=wT$, so $w_t=vw\,T\cdot AT$ and
$T_t=v(AT-(T\cdot AT)T)$. For $\rho>\varepsilon$, $A=\Pi/\rho$ and $T=\cos\phi\,\rhat-\sin\phi\,n$ ($n=\rhat^\perp$), so
$\Pi T=-\sin\phi\,n$, $T\cdot AT=\sin^2\!\phi/\rho$, and
\[
AT-(T\cdot AT)T=-\tfrac1\rho[\sin\phi\,n+\sin^2\!\phi\,T]=-\tfrac1\rho[\sin\phi\cos^2\!\phi\,n+\sin^2\!\phi\cos\phi\,\rhat]
=-\frac{\sin\phi\cos\phi}{\rho}\,T^\perp,
\]
because $T^\perp=\cos\phi\,n+\sin\phi\,\rhat$. So $\psi_t=-(v/\rho)\sin\phi\cos\phi$.
(iii) $\partial_s\rho=T\cdot\rhat=\cos\phi$. $\partial_s\rhat=\Pi T/\rho=-(\sin\phi/\rho)n$, so
$\partial_s\arg\rhat=-\sin\phi/\rho$, and $\phi=\arg\rhat-\psi$ gives $\partial_s\phi=-\sin\phi/\rho-k$. Since
$k=\partial_\theta\psi/w$, $k_t=\partial_s(\psi_t)-k\,w_t/w$, and
\[
\partial_s\Bigl(-\frac v{2\rho}\sin2\phi\Bigr)=\frac v{2\rho^2}\cos\phi\sin2\phi-\frac v\rho\cos2\phi\,\partial_s\phi
=\frac v{\rho^2}\sin\phi\cos^2\!\phi+\frac v\rho\cos2\phi\Bigl(k+\frac{\sin\phi}\rho\Bigr).
\]
So $k_t=(v/\rho)k(\cos2\phi-\sin^2\!\phi)+(v/\rho^2)\sin\phi(\cos^2\!\phi+\cos2\phi)$, which is the claim. ($\partial_t$ and
$\partial_\theta$ commute on $\psi$ because $x$ is $C^2$ in $\theta$ with $\partial_t\partial_\theta^2x$ continuous, by the second
variational equation.)
(iv) $\rho\sin\phi=T\times r$ and $\partial_s(T\times r)=k\,T^\perp\times r+T\times T=-k\,T\cdot r=-k\rho\cos\phi$; and
$T\times r\to0$ as $s\to0$.
\end{proof}

\begin{lemma}[Two-point identities]\label{lem:B-2}
For labels $a<b$ and every $t$ (primes are $t$-derivatives):
\begin{enumerate}
\item[(i)] $l'=v\int_a^bT\cdot AT\,ds$ and $c'=v\,e\cdot(F(r_b)-F(r_a))=v\int_a^be\cdot AT\,ds$;
\item[(ii)] with $\tau:=T\cdot e^\perp$ and $q:=e^\perp\cdot Ae$ (where $A=A(r(s))$),
\[
E'=v\int_a^b(T-e)\cdot AT\,ds=v\int_a^b(T-e)\cdot A(T-e)\,ds-\frac v2\int_a^b(e\cdot Ae)|T-e|^2ds+v\int_a^b\tau q\,ds;
\]
\item[(iii)] $\int_a^b|T-e|^2ds=2E=N^2$, $\int_a^b\tau\,ds=0$, and $|\tau|\le|T-e|$;
\item[(iv)] if $\rho_a,\rho_b\ge\varepsilon$ and $P$ is differentiable at $t$, then $X:=r_a\times r_b=r_a\times(x_b-x_a)$ satisfies
$X'=vX(1/\rho_a+1/\rho_b)-P'\times(x_b-x_a)$.
\end{enumerate}
\end{lemma}
\begin{proof}
(i) $l=\int|\partial_\theta x|d\theta$ and Lemma~\ref{lem:B-1}(i) (dominated convergence: $|w_t|\le(2v/\varepsilon)w$).
$c'=e\cdot(x_b'-x_a')$ and $F(r_b)-F(r_a)=\int_a^b\partial_sF(r(s))\,ds=\int_a^bAT\,ds$. (ii) Subtract; then
$(T-e)\cdot AT=(T-e)\cdot A(T-e)+(T-e)\cdot Ae$ and
$(T-e)\cdot Ae=[(T-e)\cdot e](e\cdot Ae)+[(T-e)\cdot e^\perp](e^\perp\cdot Ae)$, with $(T-e)\cdot e=T\cdot e-1=-|T-e|^2/2$ and
$(T-e)\cdot e^\perp=\tau$. (iii) $|T-e|^2=2-2T\cdot e$ and $\int_a^bT\,ds=x_b-x_a=ce$; $\int_a^b\tau\,ds=(x_b-x_a)\cdot e^\perp=0$.
(iv) $r'=v\rhat-P'$ for $\rho\ge\varepsilon$, so
$X'=v(\rhat_a\times r_b+r_a\times\rhat_b)-P'\times r_b-r_a\times P'=vX(1/\rho_a+1/\rho_b)-P'\times(r_b-r_a)$.
\end{proof}
The variational equation contains no $P'$. Hence $l$, $c$, $E$ are $C^1$ in $t$ for fixed labels and do not jump when $P$ turns.

\subsection{Excess calculus}

\begin{lemma}[Differential inequalities for the excess]\label{lem:B-3}
Fix labels $a<b$ and a time $t$; $l,c,e,E,N,\tau,q$ as above.
\begin{enumerate}
\item[(a)] For every constant $\bar q$: $E'\le v\int_a^ba^*(\rho)|T-e|^2ds+v\int_a^b\tau(q-\bar q)\,ds$.
\item[(b)] $E'\le vN\|AT\|_{L^2((a,b),ds)}$.
\item[(c)] \emph{(Far inequality.)} Let $K$ be the convex hull of the arc $\{x_s:a\le s\le b\}$, $\rho_m:=\operatorname{dist}(P,K)$,
and assume $\rho_m\ge\varepsilon$. Put $p_a:=r_a\times e$, so that $|p_a|$ is the distance from $P$ to the line through $x_a,x_b$. Then
\begin{align*}
E'&\le\frac v{\rho_m}\Bigl(1+\frac l{\rho_m}+\frac{2l^2}{\rho_m^2}\Bigr)N^2+\frac{2v\,l^{3/2}|p_a|N}{\rho_m^3},\\
l'&\le\frac{2v}{\rho_m}\Bigl[\frac{l\,(|p_a|+l^{1/2}N)^2}{\rho_m^2}+N^2\Bigr].
\end{align*}
\item[(d)] \emph{(Metric inequality.)} On a sub-arc $(s_1,s_2)$ with $\rho\ge\varepsilon$, let
$G(s):=l(s_1,s)+\rho_{s_1}-\rho_s=\int_{s_1}^s(1-\cos\phi)\,ds\ge0$. Then
\[
\int_{s_1}^{s_2}T\cdot AT\,ds=\int_{s_1}^{s_2}\frac{\sin^2\!\phi}\rho\,ds\le2\int_{s_1}^{s_2}\frac{dG}\rho,\qquad
\int_{s_1}^{s_2}\frac{dG}\rho\le\frac{G(s_2)}{\rho_{s_2}}+\int_{s_1}^{s_2}\frac G{\rho^2}\,ds.
\]
\item[(e)] \emph{(Duhamel form of (c).)} If on $[t_1,t_2]$ the hull stays at distance $\rho_m(t)\ge\max(\varepsilon,4l(t))$ from $P$,
then with $\alpha:=1.375\,v/\rho_m$ and $\beta:=2v\,l^{3/2}|p_a|\rho_m^{-3}$,
\[
N(t_2)\le\exp\Bigl(\int_{t_1}^{t_2}\alpha\Bigr)N(t_1)+\int_{t_1}^{t_2}\exp\Bigl(\int_t^{t_2}\alpha\Bigr)\beta(t)\,dt.
\]
\end{enumerate}
\end{lemma}
\begin{proof}
(a) In Lemma~\ref{lem:B-2}(ii), $(T-e)\cdot A(T-e)\le a^*|T-e|^2$ by Lemma~\ref{lem:B-V}(ii); the second term is $\le0$ since
$A\ge0$; and $\int\tau q=\int\tau(q-\bar q)$ since $\int\tau=0$. (b) Cauchy--Schwarz in $E'=v\int(T-e)\cdot AT$ and
Lemma~\ref{lem:B-2}(iii).

(c) On $K-P$ we have $|y|\ge\rho_m\ge\varepsilon$, so $q(s)=Q(r(s))$ with
$Q(y):=e^\perp\cdot\Pi(y)e/|y|=-(y\cdot e)(y\cdot e^\perp)/|y|^3$. In the coordinates $m:=y\cdot e$, $p:=y\cdot e^\perp$:
$\partial_pQ=-m(m^2-2p^2)/|y|^5$ and $\partial_mQ=-p(p^2-2m^2)/|y|^5$, hence $|\partial_pQ|\le1/|y|^2$ (the maximum of
$|c(3c^2-2)|$ on $[-1,1]$ is $1$) and $|\partial_mQ|\le2|p|/|y|^3$. For $s\in[a,b]$ the segment
$y(\sigma)=r_a+\sigma(x_s-x_a)$, $\sigma\in[0,1]$, lies in $K-P$; along it $dp=\delta(s)\,d\sigma$ and
$dm=((x_s-x_a)\cdot e)\,d\sigma$, where $\delta(s):=(x_s-x_a)\cdot e^\perp=\int_a^s\tau$ satisfies
$|\delta(s)|\le\int_a^b|T-e|\le l^{1/2}N$; moreover $|(x_s-x_a)\cdot e|\le l$ and $|p|\le|p_a|+|\delta(s)|$. Therefore
\[
|q(s)-q(a)|\le\frac{l^{1/2}N}{\rho_m^2}+\frac{2l\,(|p_a|+l^{1/2}N)}{\rho_m^3}.
\]
Use (a) with $\bar q=q(a)$, $a^*\le1/\rho_m$ on the arc, $\int|T-e|^2=N^2$ and $\int|\tau|\le l^{1/2}N$:
\[
E'\le\frac v{\rho_m}N^2+v\,l^{1/2}N\Bigl[\frac{l^{1/2}N}{\rho_m^2}+\frac{2l\,(|p_a|+l^{1/2}N)}{\rho_m^3}\Bigr],
\]
which is the first claim. For the length: $T\cdot AT=\sin^2\!\phi/\rho$, and $\rho\sin\phi=T\times r=e\times r+(T-e)\times r$, so
$|\rho\sin\phi|\le|p(s)|+\rho|T-e|$ with $|p(s)|=|r(s)\times e|\le|p_a|+l^{1/2}N$; hence
$\sin^2\!\phi\le2(|p_a|+l^{1/2}N)^2/\rho^2+2|T-e|^2$, and we integrate using $\rho\ge\rho_m$.

(d) $\sin^2\!\phi=(1-\cos\phi)(1+\cos\phi)\le2(1-\cos\phi)$, and $\partial_s\rho=\cos\phi$ gives $G'=1-\cos\phi$. Integrating by
parts, $\int G'/\rho=[G/\rho]+\int G\cos\phi/\rho^2\,ds\le G(s_2)/\rho_{s_2}+\int G/\rho^2$, as $G\ge0$ and $G(s_1)=0$.

(e) Let $N_\eta:=(N^2+\eta^2)^{1/2}$. By (c), $N_\eta N_\eta'=E'\le\alpha N^2+\beta N\le\alpha N_\eta^2+\beta N_\eta$ with
$\alpha=(v/\rho_m)(1+1/4+2/16)$; so $N_\eta'\le\alpha N_\eta+\beta$. Variation of constants, then $\eta\to0$.
\end{proof}

\begin{lemma}[Static geometry of the front]\label{lem:B-4}
At a fixed time:
\begin{enumerate}
\item[(a)] \emph{(sub-chord tilt)} for $a\le a'<b'\le b$:
$|e_{a'b'}-e_{ab}|\le\bigl[l(a',b')^{1/2}N(a,b)+E(a',b')\bigr]/c(a',b')$;
\item[(b)] \emph{(excess seen from a point of a chord)} let $a_0<a_1\le s_1<s$ and let $P$ lie on the segment $[x_{a_0},x_{a_1}]$.
Then $G:=l(s_1,s)+\rho_{s_1}-\rho_s\le E(a_0,s)$.
\end{enumerate}
\end{lemma}
\begin{proof}
(a) $c(a',b')e_{a'b'}-l(a',b')e_{ab}=\int_{a'}^{b'}(T-e_{ab})\,ds$ has norm
$\le l(a',b')^{1/2}(\int_a^b|T-e_{ab}|^2)^{1/2}=l(a',b')^{1/2}N(a,b)$, and the left side equals
$c(a',b')(e_{a'b'}-e_{ab})-E(a',b')e_{ab}$.
(b) $\rho_{s_1}\le|P-x_{a_1}|+l(a_1,s_1)$, $\rho_s\ge c(a_0,s)-|P-x_{a_0}|$, and
$|P-x_{a_0}|+|P-x_{a_1}|=c(a_0,a_1)\le l(a_0,a_1)$. Add.
\end{proof}

\subsection{Construction C and the reduction theorem}

Let $t_\varepsilon:=(1-\varepsilon)/(1-v)$, $\varepsilon_1:=e^{2\mu}\varepsilon$ and $\Lambda_0:=2\pi e^{2\mu}/(1-v)$.

\medskip\noindent\textbf{Construction C} (virtual greedy chase with step $\varepsilon$).
\begin{itemize}
\item \emph{Phase 0.} $P(s):=s$ (real axis) for $0\le s\le t_\varepsilon$.
\item \emph{Leg 0.} $P$ continues along the real axis at unit speed until it meets $x_0$ (label $\theta=0$), at time $t_0$. Put
$a_0:=0$.
\item \emph{Leg $j+1$} ($j\ge0$). At time $t_j$, $P(t_j)=x_{a_j}(t_j)$. Let
$a_{j+1}:=\max\{\theta\in[a_j,2\pi]:l_{t_j}(a_j,\theta)\le\varepsilon\}$. If $a_{j+1}=2\pi$, stop ($P$ is constant afterwards).
Otherwise $l_{t_j}(a_j,a_{j+1})=\varepsilon$, and $P$ runs at unit speed along the ray from $P(t_j)$ through
$x_{a_{j+1}}(t_j)$ until it meets $x_{a_{j+1}}$, at time $t_{j+1}$. (The target is distinct from $P(t_j)$: here
$a_j<a_{j+1}<2\pi$, so $c(a_j,a_{j+1};t_j)>0$; the only coincidence $x_{2\pi}=x_0$ occurs when the construction stops.)
\end{itemize}
Each step is computed from $P$ on earlier times, so $P$ depends only on $(v,\varepsilon)$ and is fixed in advance.

\begin{lemma}[Construction C]\label{lem:B-C}
\begin{enumerate}
\item[(0)] \emph{No accumulation.} Either the construction stops or $t_j\to\infty$; $P$ is a well-defined $1$-Lipschitz path on
$[0,\infty)$.
\item[(1)] On $[0,t_\varepsilon]$ the virtual candidates are the true ones, and $\rho_0(t_\varepsilon)=\varepsilon$. Leg $0$ lasts at
most $\varepsilon/(1-v)$, and $t_0\le1/(1-v)$.
\item[(2)] During leg $j+1$, with $u$ the unit vector from $P(t_j)$ to $x_{a_{j+1}}(t_j)$: $x_{a_{j+1}}(t)=P(t)+\rho(t)u$ with
$\rho'=-1+vf(\rho)\le-(1-v)$, and $x_{a_j}(t)=P(t)-\sigma(t)u$ with $\sigma\ge0$, $\sigma(t_j)=0$. So $x_{a_j}(t)$, $P(t)$,
$x_{a_{j+1}}(t)$ are collinear in this order, and the leg lasts at most $c(a_j,a_{j+1};t_j)/(1-v)\le\varepsilon/(1-v)$.
\item[(3)] Every $\theta\in[a_j,a_{j+1}]$ is seen by time $t_j$. If the construction stops, every direction is seen by the stopping
time.
\item[(4)] For $t$ in leg $j+1$: $l_t(a_j,a_{j+1})\le\varepsilon_1$, hence
$|P(t)-x_{a_j}(t)|\le c(a_j,a_{j+1};t)\le\varepsilon_1$; and $l_t(a,b)\le e^{2\mu}l_{t_j}(a,b)$ for all $a<b$.
\item[(5)] $R_0:=l_{t_0}(0,2\pi)\le\Lambda_0$.
\end{enumerate}
\end{lemma}
\begin{proof}
(0) On $[0,T]$, $|\partial_\theta x|\le e^{2vT/\varepsilon}$ by Lemma~\ref{lem:B-1}(i) and Lemma~\ref{lem:B-V}(ii). A non-final leg
$j+1$ with $t_j\le T$ has $\varepsilon=l_{t_j}(a_j,a_{j+1})\le e^{2vT/\varepsilon}(a_{j+1}-a_j)$, so at most
$2\pi e^{2vT/\varepsilon}/\varepsilon$ legs start before $T$.
(1) A true candidate has $\rho_\theta'=v-\rhat\cdot P'\ge-(1-v)$, so $\rho_\theta(t)\ge1-(1-v)t\ge\varepsilon$ on $[0,t_\varepsilon]$;
there the true candidates solve the virtual equation, so they are the virtual ones by uniqueness. Candidate $0$ stays on the positive
axis with $\rho_0=1-(1-v)t$. On leg $0$ the ansatz $x_0=P+\rho u$, $u=1$, solves both equations with
$\rho'=-1+vf(\rho)\le-(1-v)$, $\rho(t_\varepsilon)=\varepsilon$.
(2) The same ansatz for the target. For $a_j$: the ansatz $x_{a_j}=P-\sigma u$ (then $r=-\sigma u$, $F(r)=-f(\sigma)u$, $P'=u$)
solves $x'=vF(r)$ if and only if $\sigma'=1+vf(\sigma)$, $\sigma(t_j)=0$, so $\sigma\ge0$; conclude by uniqueness. Finally
$c(a_j,a_{j+1};t_j)\le l_{t_j}(a_j,a_{j+1})\le\varepsilon$, and $P$ is piecewise linear with speed $\le1$.
(3) $\rho_\theta(t_j)=c(a_j,\theta;t_j)\le l_{t_j}(a_j,\theta)\le\varepsilon$; apply Lemma~\ref{lem:B-VC}(c). The intervals
$[a_j,a_{j+1}]$ cover $[0,2\pi]$.
(4) $l'=v\int T\cdot AT\le(2v/\varepsilon)l$ by Lemma~\ref{lem:B-2}(i) and Lemma~\ref{lem:B-V}(ii), and the leg lasts at most
$\varepsilon/(1-v)$; the middle claim follows from (2).
(5) By (1) and Proposition~\ref{prop:length}(2), $l_{t_\varepsilon}(0,2\pi)\le2\pi+2\pi vt_\varepsilon<2\pi/(1-v)$; leg $0$
multiplies lengths by at most $e^{2\mu}$.
\end{proof}

\medskip\noindent\textbf{Hypotheses.} Fix constants $\gamma\in(0,1]$, $r_0\in(0,1]$, $c_0,c_1\ge0$ with $c_0+c_1\le1/4$,
$r_1\in(0,1]$, $K_T\ge0$, a label $b^*\in(0,2\pi]$, a starting leg index $\Bjs\ge0$, and a bound
$\BLs\ge R_{\Bjs}:=l_{t_{\Bjs}}(a_{\Bjs},2\pi)$ (for $\Bjs=0$ one may take $\BLs=\Lambda_0$). The following are required for every
leg $j+1>\Bjs$ reached by the construction and every $t\in[t_j,t_{j+1}]$.
\begin{itemize}
\item[\textup{(F$_b$)}] \emph{Flatness of the bulk seen from the base.} For every label $s\in(a_{j+1},b^*]$ with
$l:=l_t(a_j,s)\le r_0$:
\[
E(a_j,s;t)\le\Phi(l):=c_0\,l\,(l/r_0)^{2\gamma}+c_1\varepsilon .
\]
\item[\textup{(S)}] \emph{Separation.} For every $s\in(a_{j+1},b^*]$ with $l_t(a_j,s)>r_0$: $\rho_s(t)\ge r_1$.
\item[\textup{(T)}] \emph{Tip.} $K_{\mathrm{tip}}(t):=\displaystyle\int_{\max(a_{j+1},b^*)}^{2\pi}T\cdot AT\,ds\le K_T$.
\end{itemize}

\begin{theorem}[R: reduction]\label{thm:R}
Assume \textup{(F$_b$)}, \textup{(S)}, \textup{(T)} and put
\[
K^*:=16e^{2\mu}+4c_0(1+1/\gamma)+2c_1+2e^{2\mu}\BLs/r_1+K_T .
\]
If $\mu K^*<1$, construction C stops after at most $\BLs/(\varepsilon(1-\mu K^*))$ legs after leg $\Bjs$, at a time
\[
T\le t_{\Bjs}+\frac{\BLs}{(1-v)(1-\mu K^*)}\qquad\Bigl(\text{for }\Bjs=0:\ T\le\frac1{1-v}+\frac{\Lambda_0}{(1-v)(1-\mu K^*)}\Bigr),
\]
and every direction is seen by then.
\end{theorem}
\begin{proof}
Let $R_j:=l_{t_j}(a_j,2\pi)$ and, for $t$ in leg $j+1$, $\tilde K(t):=\int_{a_{j+1}}^{2\pi}T\cdot AT\,ds$.

\emph{Claim: if $R_j\le\BLs$ then $\tilde K\le K^*$ on leg $j+1$.} The labels $>b^*$ contribute at most $K_T$ by \textup{(T)}. The
labels $s\in(a_{j+1},b^*]$ are split according to $l:=l_t(a_j,s)$, which is strictly increasing in $s$ with $ds=dl$.

\emph{Near part, $l\le8\varepsilon_1$.} Its measure is $\le8\varepsilon_1$ and $T\cdot AT\le a^*\le2/\varepsilon$: contribution
$\le16e^{2\mu}$.

\emph{Far part, $l>r_0$.} $\rho\ge r_1$ by \textup{(S)}, so $T\cdot AT\le a^*(\rho)\le2/r_1$ (if $\rho\ge\varepsilon$:
$1/\rho\le1/r_1$; if $\rho<\varepsilon$: $2/\varepsilon<2/\rho\le2/r_1$). The length is
$\le l_t(a_{j+1},2\pi)\le e^{2\mu}R_j\le e^{2\mu}\BLs$ by Lemma~\ref{lem:B-C}(4): contribution $\le2e^{2\mu}\BLs/r_1$.

\emph{Middle part, $8\varepsilon_1<l\le r_0$}, i.e.\ labels $s\in(s_1,s_2]$, where $s_1$ is the label with $l=8\varepsilon_1$
($s_1>a_{j+1}$ since $l_t(a_j,a_{j+1})\le\varepsilon_1$) and $l_2\le r_0$ is the value of $l$ at $s_2$. For $l\ge\varepsilon$,
$\Phi(l)\le(c_0+c_1)l\le l/4$, so $c(a_j,s)\ge3l/4$ and, by Lemma~\ref{lem:B-C}(4),
\[
\rho_s\ge c(a_j,s)-|P-x_{a_j}|\ge3l/4-\varepsilon_1\ge l/2\quad(\ge4\varepsilon_1>\varepsilon).
\]
So Lemma~\ref{lem:B-3}(d) applies on $(s_1,s_2)$, and by Lemma~\ref{lem:B-4}(b) (with $P\in[x_{a_j},x_{a_{j+1}}]$,
Lemma~\ref{lem:B-C}(2)) $G(s)\le E(a_j,s)\le\Phi(l)$. Hence the contribution is at most
\[
2\Bigl[\frac{\Phi(l_2)}{l_2/2}+\int_{8\varepsilon_1}^{l_2}\frac{4\Phi(l)}{l^2}\,dl\Bigr]
\le4\Bigl(c_0+\frac{c_1}8\Bigr)+8\Bigl[\frac{c_0}{2\gamma}+\frac{c_1\varepsilon}{8\varepsilon_1}\Bigr]\le4c_0(1+1/\gamma)+2c_1 .
\]
This proves the claim.

\emph{Induction.} $R_{\Bjs}\le\BLs$ (for $\Bjs=0$ by Lemma~\ref{lem:B-C}(5)). If $R_j\le\BLs$ ($j\ge\Bjs$) and leg $j+1$ is not
final, then by Lemma~\ref{lem:B-2}(i)
\[
R_{j+1}=l_{t_j}(a_{j+1},2\pi)+v\int_{t_j}^{t_{j+1}}\tilde K\,dt\le R_j-\varepsilon+\frac{vK^*\varepsilon}{1-v}
=R_j-\varepsilon(1-\mu K^*)\le\BLs .
\]
So $R_j\le\BLs-(j-\Bjs)\varepsilon(1-\mu K^*)$; as $R_j\ge0$ there are at most $\BLs/(\varepsilon(1-\mu K^*))$ non-final legs after
$\Bjs$, each of duration $\le\varepsilon/(1-v)$ (and $t_0\le1/(1-v)$). Everything is seen by Lemma~\ref{lem:B-C}(3).
\end{proof}

\begin{remark}\label{rem:B-R}
(1) Without any hypothesis the proof gives a criterion: if $\tilde K(t)\le K^*$ for all $t\ge t_{\Bjs}$ and $\mu K^*<1$, the same
conclusion holds ($\tilde K$ is the quantity $K$ of Proposition~\ref{prop:D6} for the virtual front beyond the current target).
(2) The proof uses the path only from $t_{\Bjs}$ on and the bound $\BLs$. What Harry did before $t_{\Bjs}$ is irrelevant, as long as
all directions with labels $<a_{\Bjs}$ have been seen.
\end{remark}

\subsection{The bulk theorem}

Standing assumptions in this subsection: construction C and $v\le10^{-12}$, so that $e^{2\mu}\le2$, $\varepsilon_1\le2\varepsilon$
and every leg lasts at most $2\varepsilon$.

\emph{Labels.} Fix a leg start $\Bts:=t_{\Bjs}$ and a label $b^*\in(a_{\Bjs},2\pi]$. Re-label the arc $[a_{\Bjs},b^*]$ of the front
by its arclength at time $\Bts$ measured from $x_{a_{\Bjs}}$: $u\in[0,U]$, $U:=l_{\Bts}(a_{\Bjs},b^*)$. For labels $a<b$ in
$[0,U]$ let $\lambda:=u_b-u_a=l_{\Bts}(a,b)$ be the \emph{label length}. Lengths and chords are non-decreasing in $t$
(Lemma~\ref{lem:B-VC}), so $l_t(a,b)\ge\lambda$ and $c(a,b;t)\ge c(a,b;\Bts)$ for $t\ge\Bts$.

\medskip\noindent\textbf{Start hypothesis (ST)} at $\Bts$, with constants $\kappa>0$ and $R\in(0,1]$:
\begin{itemize}
\item[(IC1)] $E(a,b;\Bts)\le\kappa^2\lambda^3$ for all pairs in $[0,U]$ with $\varepsilon/2\le\lambda\le32R$;
\item[(IC2)] $c(a,b;\Bts)\ge16R$ for all pairs in $[0,U]$ with $\lambda>32R$;
\item[(IC3)] $\kappa R\le1/1584$ and $\varepsilon\le R/64$.
\end{itemize}
By (IC1), (IC3), for $\varepsilon/2\le\lambda\le32R$ we have $E\le\lambda(32\kappa R)^2\le\lambda/2$, so $c(a,b;\Bts)\ge\lambda/2$.
Hence, for all pairs with $\lambda\ge\varepsilon/2$ and all $t\ge\Bts$,
\begin{equation}\label{eq:B-CA}
c(a,b;t)\ge\tfrac12\min(\lambda,32R).\tag{CA}
\end{equation}

\emph{Family.} Let $t$ be in leg $j+1$ ($t_j\le t\le t_{j+1}$, $j\ge\Bjs$) and assume $u_{a_{j+1}}\le U$, i.e.\ the target lies
inside the labelled arc. The family $\BF(t)$ consists of the pairs $(a,b)$ with $u_b\le U$, $\varepsilon/2\le\lambda\le R$, and
($a=a_j$ or $a\ge a_{j+1}$). The \emph{gap} of such a pair is $g:=u_a-u_{a_j}\ge0$. A pair is \emph{far} if $g\ge64\lambda$ and
\emph{near} otherwise. The family only shrinks as $t$ grows, and $g$ is non-increasing in $t$.

\begin{theorem}[B: bulk]\label{thm:Bulk}
Assume \textup{(ST)}, $v\le10^{-12}$ and $v(4U+2)\le10^{-6}R$. Put $M:=R^{-1/4}/140$. Then for every $t\ge\Bts$ up to the end of the
last leg $j+1$ with $u_{a_{j+1}}\le U$, and every pair in $\BF(t)$:
\begin{equation}\label{eq:B-I}
\textup{(I1)}\quad l_t(a,b)\le2\lambda;\qquad\qquad\textup{(I2)}\quad E(a,b;t)\le\tfrac12M^2\lambda^{3/2}\quad(\text{i.e.\ }N\le M\lambda^{3/4}).
\end{equation}
Consequently, if $b'$ is a label with $u_{b'}\le U-\varepsilon/2$, then for every leg $j+1>\Bjs$ with $a_{j+1}<b'$ and every $t$ in it,
hypotheses \textup{(F$_b$)} and \textup{(S)} of Theorem~\ref{thm:R} hold with $b'$ in the place of $b^*$ and
\[
\gamma=\tfrac14,\qquad r_0=R,\qquad c_0=\tfrac12M^2R^{1/2}=\tfrac1{39200},\qquad c_1=0,\qquad r_1=\tfrac R8 .
\]
\end{theorem}

No constant depends on $\varepsilon$; $\varepsilon$ enters only through $\varepsilon\le R/64$ and $\lambda\ge\varepsilon/2$. The proof
is a continuity argument on \eqref{eq:B-I}; its five steps occupy the rest of this subsection.

\begin{proof}[Proof of the consequence, given \eqref{eq:B-I}]
Here $u_{a_{j+1}}<u_{b'}\le U-\varepsilon/2$, so the leg is one of those considered, and Step~1(a) below gives
$u_{a_{j+1}}-u_{a_j}\ge\varepsilon/2$. Let $s\in(a_{j+1},b']$; the pair $(a_j,s)$ has $\lambda\ge\varepsilon/2$. If
$l:=l_t(a_j,s)\le R$, then $\lambda\le l\le R$, the pair is in $\BF(t)$, and
$E\le\frac12M^2\lambda^{3/2}\le\frac12M^2R^{1/2}\,l\,(l/R)^{1/2}$: this is \textup{(F$_b$)}. If $l>R$, then $\lambda>R/2$ (if
$\lambda\le R$, (I1) gives $\lambda\ge l/2$), and by \eqref{eq:B-CA} and Lemma~\ref{lem:B-C}(4)
$\rho_s\ge c(a_j,s;t)-|P-x_{a_j}|\ge\frac12\min(\lambda,32R)-2\varepsilon\ge R/4-R/32\ge R/8$: this is \textup{(S)}.
\end{proof}

\begin{proof}[Proof of \eqref{eq:B-I}: the continuity argument]
Let $t^*$ be the supremum of the times $t'$ such that (I1), (I2) hold for all $t\in[\Bts,t']$ and all pairs of $\BF(t)$. At $\Bts$:
$l=\lambda$ and $E\le\kappa^2\lambda^3\le\kappa^2R^{3/2}\lambda^{3/2}\le M^2\lambda^{3/2}/128$, since (IC3) is
$\kappa R^{3/4}\le M/(8\sqrt2)$. We shall show:
\begin{quote}
\textbf{(Imp)} if (I1), (I2) hold on $[\Bts,t^*)$, then for every $t\le t^*$ and every pair of $\BF(t)$:
$l_t\le\sqrt2\,\lambda$ and $E\le\frac14M^2\lambda^{3/2}$.
\end{quote}
This proves the theorem. Indeed the bounds of (Imp) hold at $t^*$ by continuity (a pair of $\BF(t^*)$ is in $\BF(t)$ for $t<t^*$).
For $t\in[t^*,t^*+\delta]$: $l_t\le e^{2v\delta/\varepsilon}l_{t^*}$ (as $l'\le(2v/\varepsilon)l$) and
$|E'|\le v\int|T-e||AT|\,ds\le(4v/\varepsilon)l$, so $E(t)\le\frac14M^2\lambda^{3/2}+\delta(8v/\varepsilon)\cdot2\lambda$. Both stay
below the bounds (I1), (I2) for a $\delta>0$ that depends only on $(v,\varepsilon,M)$, because $\lambda\ge\varepsilon/2$. Since
$\BF(t)\subset\BF(t^*)$, $t^*$ would not be the supremum unless it is the final time. (By Lemma~\ref{lem:B-C}(0) there are finitely
many legs on bounded time intervals.)

From now on (I1), (I2) are assumed on $[\Bts,t^*)$, all times are $<t^*$, and ``family bounds'' refers to them. Steps~1--5 below prove
(Imp).
\end{proof}

\paragraph{Step 1 (legs, times, distances).}
(a) \emph{For the legs considered ($u_{a_{j+1}}\le U$): $u_{a_{j+1}}-u_{a_j}\in[\varepsilon/2,\varepsilon]$.} It is
$\le l_{t_j}(a_j,a_{j+1})\le\varepsilon$. If it were $<\varepsilon/2$, the label $s$ with $u_s=u_{a_j}+\varepsilon/2$ (which is
$\le U$, since otherwise $a_{j+1}$ would have label $>U$) would have $l_{t_j}(a_j,s)\le2\cdot\varepsilon/2=\varepsilon$ by (I1) and
$s>a_{j+1}$, contradicting the maximality of $a_{j+1}$.

(b) Hence at most $2U/\varepsilon+1$ legs are considered, and the time elapsed since $\Bts$ is at most $T_B:=4U+2$.

(c) For a fixed label $a$ the gaps at successive legs are spaced by $\ge\varepsilon/2$, and each leg lasts $\le2\varepsilon$. So for
every non-increasing $h\ge0$ and every $g_0\ge0$,
\begin{equation}\label{eq:B-1c}
\int_{\{t:\,g(t)\ge g_0\}}h(g(t))\,dt\le2\varepsilon\,h(g_{\min})+4\int_{g_{\min}}^\infty h(g)\,dg,
\end{equation}
where $g_{\min}$ is the smallest gap $\ge g_0$ attained. (For every attained gap $g_i>g_{\min}$,
$h(g_i)\le(2/\varepsilon)\int_{g_i-\varepsilon/2}^{g_i}h$, and these intervals are disjoint in $[g_{\min},\infty)$.)

(d) For $t$ in leg $j+1$ and a label $s$ with $w_s:=u_s-u_{a_j}\ge0$, by \eqref{eq:B-CA} and Lemma~\ref{lem:B-C}(4),
\begin{equation}\label{eq:B-star}
\rho_s\ge\tfrac12\min(w_s,32R)-2\varepsilon,\qquad\text{hence}\qquad\max(\rho_s,\varepsilon)\ge\tfrac16\min(w_s,32R).\tag{$*$}
\end{equation}
(With $m:=\min(w_s,32R)$: if $m\le6\varepsilon$ then $\varepsilon\ge m/6$; if $m>6\varepsilon$ then $m/2-2\varepsilon\ge m/6$. For
$w_s<\varepsilon/2$, where \eqref{eq:B-CA} is not available, the first inequality is trivial.)

\paragraph{Step 2 (tilt lemma).}
\emph{At any time, if $J\subset I$ are pairs of the family, then}
\begin{equation}\label{eq:B-TL}
|e_J-e_I|\le91\,M\lambda_I^{1/4}.\tag{TL}
\end{equation}
First let $J_1\subset I_1$ be family pairs with $\lambda_{I_1}\le3\lambda_{J_1}$. By Lemma~\ref{lem:B-4}(a), (I1), (I2) and
\eqref{eq:B-CA},
\[
|e_{J_1}-e_{I_1}|\le\Bigl[(2\lambda_{J_1})^{1/2}M(3\lambda_{J_1})^{3/4}+\tfrac12M^2\lambda_{J_1}^{3/2}\Bigr]\frac2{\lambda_{J_1}}
=M\lambda_{J_1}^{1/4}\bigl[2\sqrt2\cdot3^{3/4}+M\lambda_{J_1}^{1/4}\bigr]\le6.7\,M\lambda_{J_1}^{1/4},
\]
as $MR^{1/4}=1/140$ and $2\sqrt2\cdot3^{3/4}+1/140=6.455$. Now let $J=(a',b')\subset I=(a,b)$.

(A) Extend to the right: $b_k:=\min(u_{a'}+2^k\lambda_J,u_b)$, $k=0,\dots,k_1$, with $b_{k_1}=b$. The arcs $(a',b_k)$ are in the family
(same first point $a'$, length $\le R$), consecutive label lengths have ratio $\le2$, and
\[
\sum_{k<k_1}6.7M(2^k\lambda_J)^{1/4}\le\frac{6.7\,M\lambda_I^{1/4}}{1-2^{-1/4}}\le42.12\,M\lambda_I^{1/4}
\qquad\bigl(6.7/(1-2^{-1/4})=42.1109\ldots\bigr).
\]

(B) If $a'>a$, extend to the left down to $a^\circ:=a$ if $a\ge a_{j+1}$, and $a^\circ:=a_{j+1}$ if $a=a_j$: put
$a_k':=\max(u_b-2^k\lambda(a',b),u_{a^\circ})$. The arcs $(a_k',b)$ start at labels $\ge a_{j+1}$, so they are in the family; the
same sum is $\le42.12\,M\lambda_I^{1/4}$.

(C) If $a=a_j\ne a^\circ$: $(a_{j+1},b)\subset(a_j,b)$ with
$\lambda(a_j,b)\le\lambda(a_{j+1},b)+\varepsilon\le3\lambda(a_{j+1},b)$ (as $\lambda(a_{j+1},b)\ge\lambda_J\ge\varepsilon/2$): one
more step, $\le6.7M\lambda_I^{1/4}$.

The total is $\le(2\cdot42.12+6.7)M\lambda_I^{1/4}=90.94\,M\lambda_I^{1/4}$, which proves \eqref{eq:B-TL}.

\emph{Consequence.} For a family pair $(a,b)$ with label length $\lambda$ and any family sub-arc $J\subset(a,b)$ with label length
$\lambda_J$ (in particular $J=(a,s)$ with $u_s-u_a\ge\varepsilon/2$; shorter arcs are always enlarged to label length
$\varepsilon/2$ where this is used),
\begin{equation}\label{eq:B-mu}
\int_J|T-e_{ab}|^2ds\le2N(J)^2+2l(J)|e_J-e_{ab}|^2\le2M^2\lambda_J^{3/2}+4\lambda_J\cdot91^2M^2\lambda^{1/2}
\le C_\mu M^2\lambda^{1/2}\lambda_J,
\end{equation}
with $C_\mu:=33126$, and by Cauchy--Schwarz
$\int_J|T-e_{ab}|\,ds\le(2\lambda_J\cdot C_\mu M^2\lambda^{1/2}\lambda_J)^{1/2}\le258\,M\lambda^{1/4}\lambda_J$.

\paragraph{Step 3 (far pairs: (I2) with $M/2$).}
Let $(a,b)$ be far at time $t$ (in leg $j+1$): $g\ge64\lambda\ge32\varepsilon$, so $a>a_{j+1}$, and the pair has been in the family
and far since $\Bts$. Put $\hat g:=\min(g,32R)$. The hull of the arc lies in $B(x_a,l)$ with $l\le2\lambda$, so by \eqref{eq:B-star}
$\rho_m\ge\hat g/2-2\varepsilon-2\lambda$. If $g\le32R$ this is $\ge g/2-6\lambda$ (as $2\varepsilon\le4\lambda$)
$\ge(\frac12-\frac6{64})g$, which is $\ge g/4$ and $\ge26\lambda$. If $g>32R$ it is $16R-2\varepsilon-2\lambda\ge13.9R$, which is
$\ge8R=\hat g/4$ and $\ge13.9\lambda$. In both cases
\[
\rho_m\ge\hat g/4\qquad\text{and}\qquad\rho_m\ge13.9\lambda\ge\max(\varepsilon,4l)
\]
($l\le2\lambda$, $\varepsilon\le2\lambda$), so Lemma~\ref{lem:B-3}(e) applies on the whole far period, with
$\alpha\le1.375v/\rho_m\le5.5v/\hat g$.

\emph{Amplification.} Between a time with gap $g_1$ and a later one with gap $g_2$ (both $\le32R$), by \eqref{eq:B-1c},
$\int\alpha\,dt\le5.5v[2\varepsilon/g_2+4\log(g_1/g_2)]\le0.35v+22v\log(g_1/g_2)$; while $g>32R$,
$\int\alpha\,dt\le5.5vT_B/(32R)$. So with $\nu:=22v$, $A_0:=\exp(v(0.35+0.18\,T_B/R))\le2$ and
$\hat g_2:=\min(g_2,32R)$, for arbitrary gaps $g_1\ge g_2$,
\[
\exp\Bigl(\int\alpha\Bigr)\le A_0\bigl(\min(g_1,32R)/\hat g_2\bigr)^\nu
\]
(if $g_2>32R$ only the second bound is used and the power is $1$).

\emph{Forcing} $\beta=2v\,l^{3/2}|p_a|/\rho_m^3$. If $g>R/2$: $|p_a|\le\rho_a\le\rho_m+l$ ($x_a$ lies in the hull, whose diameter is
$\le l$). For $R/2<g\le32R$: $l\le2\lambda\le\rho_m/13$ and $\rho_m\ge(\frac12-\frac6{64})R/2\ge R/8$, so
$\beta\le2v(2\lambda)^{3/2}\cdot\frac{14}{13}\cdot64/R^2\le390\,v\lambda^{3/2}/R^2$. For $g>32R$: $l\le2R\le0.144\rho_m$ and
$\rho_m\ge13.9R$, so $\beta\le2v(2\lambda)^{3/2}\cdot1.144/(13.9R)^2\le0.04\,v\lambda^{3/2}/R^2$. Hence
\[
\beta\le408\,v\lambda^{3/2}/R^2\qquad\text{whenever }g>R/2 .
\]
If $g\le R/2$: $P$ lies on $[x_{a_j},x_{a_{j+1}}]$ (Lemma~\ref{lem:B-C}(2)), so
$|p_a|\le\max_{a'\in\{a_j,a_{j+1}\}}\operatorname{dist}(x_{a'},\operatorname{line}(x_a,x_b))$, and this distance equals
$c(a',a)|e_{a'a}\times e_{ab}|\le c(a',a)(|e_{a'a}-e_{a'b}|+|e_{ab}-e_{a'b}|)$. The pairs $(a',a)$, $(a,b)$, $(a',b)$ are in the
family (label lengths $\ge64\lambda-\varepsilon\ge\varepsilon/2$, resp.\ $\le g+\lambda\le R$), so by \eqref{eq:B-TL} and
$c(a',a)\le l(a',a)\le2g$:
\[
|p_a|\le2g\cdot2\cdot91M(g+\lambda)^{1/4}\le366\,Mg^{5/4},\qquad
\beta\le2v(2\lambda)^{3/2}\cdot366Mg^{5/4}\cdot\frac{64}{g^3}\le1.33\cdot10^5\,vM\lambda^{3/2}g^{-7/4}.
\]

\emph{Result.} At a time with gap $g_2\ge64\lambda$, using $N(\Bts)\le\sqrt2\kappa\lambda^{3/2}$ and \eqref{eq:B-1c} with
$h(g)=g^{\nu-7/4}$, $\nu\le1/12$, and with $\hat g_2=\min(g_2,32R)$:
\begin{align*}
N&\le A_0\Bigl(\frac{32R}{\hat g_2}\Bigr)^{\!\nu}\Bigl[\sqrt2\kappa\lambda^{3/2}+\frac{408\,vT_B\lambda^{3/2}}{R^2}\Bigr]
+A_0\cdot1.33\cdot10^5\,vM\lambda^{3/2}g_2^{-\nu}\Bigl[2\varepsilon g_2^{\nu-7/4}+\frac{4g_2^{\nu-3/4}}{3/4-\nu}\Bigr]\\
&\le A_0\Bigl[\bigl(\sqrt2\kappa+408\,vT_B/R^2\bigr)R^{3/4}+3.6\cdot10^4\,vM\Bigr]\lambda^{3/4},
\end{align*}
because
\[
(32R/\hat g_2)^\nu\lambda^{3/2}\le\max\bigl(1,R/(2\lambda)\bigr)^\nu\lambda^{3/2}\le(R/\lambda)^\nu\lambda^{3/2}\le
R^{3/4}\lambda^{3/4},
\]
and moreover
$g_2^{-3/4}(1/16+6)=6.0625\,g_2^{-3/4}\le6.0625\,(64\lambda)^{-3/4}\le0.268\,\lambda^{-3/4}$. Since $\sqrt2\kappa R^{3/4}\le M/8$,
$408\,vT_BR^{-5/4}\le408\cdot10^{-6}R^{-1/4}\le M/16$ and $3.6\cdot10^4v\le1/16$, we get
\[
N\le2\cdot(M/4)\lambda^{3/4}=(M/2)\lambda^{3/4},\qquad\text{i.e.}\qquad E\le M^2\lambda^{3/2}/8\quad\text{for far pairs.}
\]

\paragraph{Step 4 (near pairs).}
Let $(a,b)\in\BF(t)$ be near ($a=a_j$, $g=0$; or $a\ge a_{j+1}$, $g<64\lambda$). The total time a given pair spends near is
$\le(128\lambda/\varepsilon+1)\cdot2\varepsilon+2\varepsilon\le264\lambda$ (Step~1(c); the last $2\varepsilon$ is the leg in which $a$
is the base). When it becomes near (or at $\Bts$), $E\le M^2\lambda^{3/2}/8$ (Step~3, resp.\ the start; $E$ is continuous at leg
switches). We bound the increase $\Delta E$ of $E$ while near by integrating Lemma~\ref{lem:B-3}(a) with $\bar q=0$:
\[
E'\le\mathrm I+\mathrm{II},\qquad\mathrm I:=v\int_a^ba^*(\rho)|T-e|^2ds,\qquad\mathrm{II}:=v\int_a^b|\tau||q|\,ds,\qquad e:=e_{ab}.
\]
Put $\bar\rho:=\max(\varepsilon,\min_{[a,b]}\rho)$ and $\tilde g:=\max(g,\varepsilon)$. Every $s\in[a,b]$ has $w_s\ge g$, so by
\eqref{eq:B-star} $\bar\rho\ge\min(g,32R)/6$. A near pair can have $g\in(32R,64R)$ when $\lambda>R/2$; then
$6\bar\rho\ge32R\ge\lambda$ and $\log^+(\lambda/(6\bar\rho))=0$; otherwise $6\bar\rho\ge\tilde g$. So
$\log^+(\lambda/(6\bar\rho))\le\log^+(\lambda/\tilde g)$ in all cases.

\emph{Term \textup{I}.} $a^*(\rho)=(1/\varepsilon)\mathbf 1_{\rho<\varepsilon}+\int_{\max(\rho,\varepsilon)}^\infty y^{-2}dy$, so
\[
\mathrm I=v\Bigl[\frac{\omega(\varepsilon)}\varepsilon+\int_\varepsilon^\infty\omega(y)y^{-2}dy\Bigr],\qquad
\omega(y):=\int_{\{s\in(a,b):\,\rho_s<y\}}|T-e|^2ds .
\]
By \eqref{eq:B-star}, $\rho_s<y$ forces $w_s<2y+4\varepsilon$ when $2y+4\varepsilon\le32R$; so these $s$ lie in the sub-arc $J_y$ of
$(a,b)$ starting at $a$ with label length $\max(\min(\lambda,2y+4\varepsilon-g),\varepsilon/2)\le\min(\lambda,6y)$ (for
$y\ge\varepsilon$), a family pair. By \eqref{eq:B-mu}, $\omega(y)\le C_\mu M^2\lambda^{1/2}\min(\lambda,6y)$ (also for larger $y$,
trivially), and $\omega(y)=0$ for $y\le\min\rho$. Hence
\[
\mathrm I\le vC_\mu M^2\lambda^{1/2}\Bigl[6+\int_{\bar\rho}^{\lambda/6}\frac{6\,dy}y+\int_{\lambda/6}^\infty\frac{\lambda\,dy}{y^2}\Bigr]
\le6vC_\mu M^2\lambda^{1/2}\bigl[2+\log^+(\lambda/\tilde g)\bigr].
\]

\emph{Term \textup{II}.} By Lemma~\ref{lem:B-V}(iii), $|q|\le|p(s)|/\max(\rho_s,\varepsilon)^2$ with $p(s):=(x_s-P)\times e$, and
since $P\in[x_{a_j},x_{a_{j+1}}]$, $|p(s)|\le\max_{a'\in\{a_j,a_{j+1}\}}|(x_s-x_{a'})\times e|$. Let
\[
\vartheta:=260\,M\lambda^{1/4}\ \text{ if }65\lambda\le R,\qquad\vartheta:=2\ \text{ otherwise}
\]
(note $258M\lambda^{1/4}\le\vartheta\le2$ in both cases), and $\lambda_s':=\max(u_s-u_a,\varepsilon/2)$.

\emph{Case $a'\le a$} (always when $a\ge a_{j+1}$; and $a'=a_j$ when $a=a_j$).
$(x_s-x_{a'})\times e=\int_a^s\tau+(x_a-x_{a'})\times e$. The first term is
$\le\int_{(a,s)}|T-e|\le258M\lambda^{1/4}\lambda_s'\le\vartheta\lambda_s'$ (Step~2). The second vanishes if $a'=a$. Assume first
$65\lambda\le R$. If $g':=u_a-u_{a'}\ge\varepsilon/2$, the second term equals
$c(a',a)|e_{a'a}\times e_{ab}|\le2g'(|e_{a'a}-e_{a'b}|+|e_{ab}-e_{a'b}|)\le4\vartheta g$: the three pairs $(a',a)$, $(a,b)$, $(a',b)$
are in the family ($\lambda(a',b)\le65\lambda\le R$), $c(a',a)\le2g'$ by (I1), and \eqref{eq:B-TL} bounds each tilt by
$91M(65\lambda)^{1/4}\le\vartheta$. If $0<g'<\varepsilon/2$ (only $a'=a_{j+1}$): with $J_0:=(a',\text{label }u_{a'}+\varepsilon/2)
\subset(a',b)$,
\[
|(x_a-x_{a'})\times e|\le l(a',a)^{1/2}\Bigl(\int_{J_0}|T-e_{ab}|^2\Bigr)^{1/2}
\le\varepsilon^{1/2}\bigl[2M^2(\varepsilon/2)^{3/2}+2\varepsilon(2\vartheta)^2\bigr]^{1/2}\le3\vartheta\varepsilon .
\]

\emph{Case $a=a_j$, $a'=a_{j+1}\le b$.} $(x_s-x_{a'})\times e=\pm\int\tau$ over the arc between $a'$ and $s$, which is contained in
$(a,\max(s,a'))\subset(a,b)$; so it is $\le\vartheta\max(\lambda_s',\varepsilon)$ by Step~2.

\emph{Case $a=a_j$, $a'=a_{j+1}>b$} (possible when $\varepsilon/2\le\lambda<u_{a_{j+1}}-u_{a_j}$; then $\lambda<\varepsilon$). Here
the arc between $s$ and $a'$ is not a sub-arc of $(a,b)$. Put $I:=(a_j,a_{j+1})$, a family pair containing $(a,b)$, with
$\varepsilon/2\le\lambda_I\le\varepsilon\le2\lambda$. For $s\in[a,b]$, by (I1), (I2) for $I$ and \eqref{eq:B-TL} applied to
$(a,b)\subset I$,
\begin{align*}
|(x_s-x_{a'})\times e|\le\int_I|T-e_{ab}|\,ds&\le\int_I|T-e_I|\,ds+l(I)\,|e_I-e_{ab}|
\le(2\lambda_I)^{1/2}M\lambda_I^{3/4}+2\lambda_I\cdot91M\lambda_I^{1/4}\\
&=(\sqrt2+182)M\lambda_I^{5/4}\le183.42\,M\varepsilon^{1/4}\varepsilon\le219\,M\lambda^{1/4}\varepsilon\le\vartheta\varepsilon
\end{align*}
($183.42\cdot2^{1/4}=218.12$).

So if $65\lambda\le R$: $|p(s)|\le\vartheta(\lambda_s'+4g+3\varepsilon)\le\vartheta(5w_s+4\varepsilon)$ (as
$\lambda_s'\le w_s-g+\varepsilon/2$), and with \eqref{eq:B-star}, for $w_s\le32R$,
\[
|q|\le\frac{\vartheta(5w_s+4\varepsilon)}{\max(\varepsilon,w_s/6)^2}\le\frac{34\vartheta}{\max(\varepsilon,w_s/6)}
=\frac{204\vartheta}{\max(w_s,6\varepsilon)}=\int_{\max(w_s,6\varepsilon)}^\infty204\vartheta\,y^{-2}dy .
\]
If $65\lambda>R$ ($\vartheta=2$) the same bound holds trivially: $|p(s)|\le\rho_s$ gives
$|q|\le1/\max(\rho_s,\varepsilon)\le6/\max(w_s,6\varepsilon)$. In both cases, for $w_s>32R$, $|p|\le\rho_s$ and $\rho_s\ge15R$ give
$|q|\le1/(15R)$. Therefore, with $\int_{\{s\in(a,b):\,w_s<y\}}|T-e|\,ds\le\vartheta\min(\lambda,y)$ for $y\ge6\varepsilon$ (Step~2, a
sub-arc of label length $\le\max(y-g,\varepsilon/2)$; zero if $y\le g$):
\[
\mathrm{II}\le204\,v\vartheta^2\int_{\max(6\varepsilon,g)}^\infty\frac{\min(\lambda,y)}{y^2}\,dy+\frac v{15R}\int_a^b|T-e|\,ds
\le204\,v\vartheta^2\bigl[1+\log^+(\lambda/\tilde g)\bigr]+\frac{0.1\,vM\lambda^{5/4}}R
\]
(using $\int_a^b|T-e|\le l^{1/2}N\le\sqrt2M\lambda^{5/4}$).

\emph{Time integrals.} By \eqref{eq:B-1c},
$\int_{\mathrm{near}}[1+\log^+(\lambda/\tilde g)]\,dt\le264\lambda+8\varepsilon\log^+(\lambda/\varepsilon)+4\int_\varepsilon^\lambda
\log(\lambda/g)\,dg\le272\lambda$ (at most two legs have $g<\varepsilon$, each contributing
$\le2\varepsilon\log^+(\lambda/\varepsilon)$; and $\varepsilon\log(\lambda/\varepsilon)\le\lambda/e$), and
$\int_{\mathrm{near}}[2+\log^+]\,dt\le536\lambda$. So
\begin{align*}
\Delta E&\le6\cdot536\,C_\mu vM^2\lambda^{3/2}+204\cdot272\,v\vartheta^2\lambda+26.4\,vM\lambda^{9/4}/R\\
&\le1.07\cdot10^8\,vM^2\lambda^{3/2}+5.6\cdot10^4\,v\vartheta^2\lambda+3700\,vM^2\lambda^{3/2},
\end{align*}
since $M\lambda^{9/4}/R\le M^2\lambda^{3/2}(\lambda/R)^{3/4}/(MR^{1/4})\le140M^2\lambda^{3/2}$. If $65\lambda\le R$:
$\vartheta^2\lambda=67600M^2\lambda^{3/2}$ and $\Delta E\le3.9\cdot10^9\,vM^2\lambda^{3/2}$. If $65\lambda>R$:
$\vartheta^2\lambda=4\lambda\le4(65/R)^{1/2}\lambda^{3/2}=4\cdot8.07\cdot19600\,M^2\lambda^{3/2}$ and
$\Delta E\le3.6\cdot10^{10}\,vM^2\lambda^{3/2}$. In both cases, as $v\le10^{-12}$,
\[
\Delta E\le M^2\lambda^{3/2}/8,\qquad\text{so}\qquad E\le M^2\lambda^{3/2}/4\quad\text{for near pairs.}
\]
With Step~3 this is the $E$-part of (Imp).

\paragraph{Step 5 ((I1) with $\sqrt2$).}
$l(\Bts)=\lambda$ and $l'=v\int T\cdot AT$.

\emph{Far, $g>R/2$:} $T\cdot AT\le1/\rho_m\le8/R$, so $(\log l)'\le8v/R$; total $\le8vT_B/R$.

\emph{Far, $g\le R/2$:} by Lemma~\ref{lem:B-3}(c), $l'\le(2v/\rho_m)[l(|p_a|+l^{1/2}N)^2/\rho_m^2+N^2]$, with
$|p_a|+l^{1/2}N\le366Mg^{5/4}+\sqrt2M\lambda^{5/4}\le367Mg^{5/4}$, $\rho_m\ge g/4$ and $N^2\le M^2\lambda^{3/2}\le lM^2g^{1/2}$. So
\[
(\log l)'\le\frac{8v}g\,(367^2\cdot16+1)M^2g^{1/2}=1.72402\cdot10^7\,vM^2g^{-1/2},
\]
and by \eqref{eq:B-1c}, $\int g^{-1/2}dt\le2\varepsilon(32\varepsilon)^{-1/2}+8(R/2)^{1/2}\le5.702\,R^{1/2}$ (using
$\varepsilon\le R/64$). The total is $\le1.72402\cdot10^7\cdot5.702\,vM^2R^{1/2}\le9.831\cdot10^7\,vM^2R^{1/2}\le5016\,v\le5200\,v$,
since $M^2R^{1/2}=1/19600$.

\emph{Near:} $T\cdot AT\le a^*(\rho)$, so $l'\le v[m(\varepsilon)/\varepsilon+\int_\varepsilon^\infty m(y)y^{-2}dy]$ with
$m(y):=|\{s\in(a,b):\rho_s<y\}|\le\min(2\lambda,12y)$ (as for $\omega(y)$, with (I1) for the sub-arc $J_y$). Hence
$l'\le12v[2+\log^+(\lambda/\tilde g)]$, and the increase while near is $\le12\cdot536\,v\lambda\le6500\,v\lambda$.

Altogether $l_t\le\lambda\exp(8vT_B/R+5200v)+6500\,v\lambda\le\sqrt2\,\lambda$. This completes the proof of (Imp), hence of
Theorem~\ref{thm:Bulk}. \qed

\begin{remark}\label{rem:B-B}
(1) The only structural requirements are: an exponent $1/2+\gamma$ with $0<\gamma<1/2$ (here $\gamma=1/4$), the smallness
$MR^\gamma\ll1$, and $v$ small against absolute constants and against $R/U$.
(2) Theorem~\ref{thm:Bulk} does not use the tip: every estimate for a pair $(a,b)$ involves only the arc from the base to $b$.
(3) No weight in the gap is needed. The amplification of far pairs is a power $(\min(g_1,32R)/g_2)^{22v}$ of a gap ratio, not
$e^{cv\cdot\#\mathrm{legs}}$; it acts on data of exponent $3/2$ and is absorbed by lowering the exponent to $3/4$. The forcing is
concentrated at gaps $g\approx\lambda$.
(4) \textup{(ST)} is monotone in $R$: if (IC1), (IC2) hold for $R$ and $\kappa R\le1/1584$, then \textup{(ST)} holds for every
$R'<R$ with $\varepsilon\le R'/64$. Indeed (IC1) then concerns fewer pairs; for $\lambda>32R$, $c\ge16R\ge16R'$; and for
$\lambda\in(32R',32R]$, $c\ge\lambda-\kappa^2\lambda^3\ge\lambda(1-(32\kappa R)^2)\ge16R'$.
\end{remark}

\begin{remark}[Relaxed start; not covered by the referee reports]\label{rem:B-relaxed}
The proof uses \textup{(ST)} at $\Bts$ only through: \eqref{eq:B-CA} at $\Bts$; the start value $E\le M^2\lambda^{3/2}/8$ for pairs
that are near at $\Bts$; and $N(\Bts)$ in the first term of the \emph{Result} of Step~3. Hence Theorem~\ref{thm:Bulk}, its
consequence and Corollary~\ref{cor:M} remain true, with the same proof and constants, if (IC1)--(IC3) are replaced by:
\begin{itemize}
\item[(W1)] $c(a,b;\Bts)\ge\frac12\min(\lambda,32R)$ for all pairs in $[0,U]$ with $\lambda\ge\varepsilon/2$;
\item[(W2)] $E(a,b;\Bts)\le\dfrac{M^2}{128}\lambda^{3/2}\min\Bigl(1,\dfrac{64\lambda}{\min(u_a,32R)}\Bigr)^{44v}$ for all pairs with
$\varepsilon/2\le\lambda\le R$ ($u_a$ is the gap at $\Bts$);
\item[(W3)] $\varepsilon\le R/64$,
\end{itemize}
where $M=R^{-1/4}/140$ and $R\in(0,1]$ is now a free parameter. (Near pairs: $E\le M^2\lambda^{3/2}/128$. Far pairs with start gap
$g_{\mathrm s}=u_a\ge64\lambda$: with $\hat g_2:=\min(g_2,32R)$,
\[
A_0(\min(g_{\mathrm s},32R)/\hat g_2)^\nu N(\Bts)\le2\cdot\frac M8\lambda^{3/4}\bigl(\min(\min(g_{\mathrm s},32R),64\lambda)/\hat g_2\bigr)^\nu
\le\frac M4\lambda^{3/4},
\]
and the two forcing terms are $\le A_0(M/16+M/16)\lambda^{3/4}$ as before, so $N\le(M/2)\lambda^{3/4}$.)
(W2) admits an additive transient near the base: a kink with excess $E_{\mathrm{kink}}(a,b)$ is allowed as long as
$\kappa^2\lambda^3+E_{\mathrm{kink}}$ satisfies (W2). For a kink of turning angle $\Theta$ at scale $\varepsilon$
($E_{\mathrm{kink}}\approx\Theta^2\varepsilon$) this requires $\Theta^2\lesssim M^2\varepsilon^{1/2}/360$, which is \emph{not} uniform as
$\varepsilon\to0$ at fixed $\Theta\sim v$. So this proof does not tolerate an $\varepsilon$-independent additive term $c_1\varepsilon$
in (I2); an $\varepsilon$-uniform treatment of the kink needs an invariant adapted to the kink tail, which is not proved.
This remark was added after the two referee reports and was not checked by either referee; the second report (\S4(2)) suggests
only the start ``$E\le M^2\lambda^{3/2}/128$ for family pairs, plus \eqref{eq:B-CA}'' (without the gap factor of (W2)), and the
first (\S7) only ``(IC1) for $\varepsilon/2\le\lambda\le32R$, plus \eqref{eq:B-CA}''. Its main claim was checked in the later audit
of the typeset text (report \texttt{audit-B.md}). The remark is not used by Theorems~\ref{thm:Main} and~\ref{thm:Main2}.
\end{remark}

\subsection{A conditional statement}

\begin{corollary}[M; conditional on \textup{(ST)} and \textup{(T)}]\label{cor:M}
Let $\kappa>0$, $R\in(0,1]$, $\BTs$, $\BLs$, $K_T\ge0$, $\delta\in(0,1]$, $v_1$ and $\varepsilon_0\le\min(\delta,R/64)$ be absolute
constants. \textbf{Suppose} that for all $\varepsilon\le\varepsilon_0$ and $v\le v_1$, construction C reaches a leg start
$\Bts=t_{\Bjs}\le\BTs$ with $a_{\Bjs}<2\pi-\delta$ and remaining length $R_{\Bjs}=l_{\Bts}(a_{\Bjs},2\pi)\le\BLs$ such that both of the
following hold.
\begin{itemize}
\item[\textup{(ST)}] \emph{Start.} On the arc $[a_{\Bjs},2\pi-\delta/2]$, labelled by its arclength $u\in[0,U]$ at time $\Bts$, with
$\lambda=u_b-u_a$:
\[
\begin{aligned}
&E(a,b;\Bts)\le\kappa^2\lambda^3&&\text{for all pairs with }\varepsilon/2\le\lambda\le32R,\\
&c(a,b;\Bts)\ge16R&&\text{for all pairs with }\lambda>32R,\\
&\kappa R\le1/1584,\quad\varepsilon\le R/64 .
\end{aligned}
\]
\item[\textup{(T)}] \emph{Tip.} For every leg $j+1>\Bjs$ and every $t$ in it,
\[
K_{\mathrm{tip}}(t)=\int_{\max(a_{j+1},\,2\pi-\delta)}^{2\pi}T\cdot AT\,ds\le K_T .
\]
\end{itemize}
The leg index $\Bjs$ (hence $a_{\Bjs}$, $\Bts$) may depend on $(v,\varepsilon)$; the constants
$\kappa,R,\BTs,\BLs,K_T,\delta$ may not. Let
\[
K^*:=32+\frac1{1960}+\frac{32\BLs}R+K_T,\qquad v_0:=\min\Bigl(v_1,\ 10^{-12},\ \frac{10^{-6}R}{4\BLs+2},\ \frac1{2K^*}\Bigr).
\]
\textbf{Then} for $v\le v_0$ and $\varepsilon\le\varepsilon_0$ every direction is seen by time $\BTs+2.1\BLs$, which does not depend
on $\varepsilon$; and for these $v$ there is a path of speed at most $1$ that catches every direction exactly, by time
$\BTs+2.1\BLs$.

The same holds if construction C is replaced by any $1$-Lipschitz path fixed in advance such that, at some time $\Bts$,
$P(\Bts)=x_a(\Bts)$ for a label $a=:a_{\Bjs}$, every label in $[0,a_{\Bjs}]$ has been seen by $\Bts$, and $P$ follows the legs of
construction C from $\Bts$ on.
\end{corollary}
\begin{proof}
$U\le R_{\Bjs}\le\BLs$, so $v(4U+2)\le10^{-6}R$ and Theorem~\ref{thm:Bulk} applies on $[a_{\Bjs},2\pi-\delta/2]$. As
$|\partial_\theta x|\ge1$, the labels satisfy $U-u_{2\pi-\delta}\ge\delta/2\ge\varepsilon/2$, so $b':=2\pi-\delta$ is admissible and
Theorem~\ref{thm:Bulk} gives \textup{(F$_b$)}, \textup{(S)} for $b^*=2\pi-\delta$ with $\gamma=1/4$, $r_0=R$, $c_0=1/39200$, $c_1=0$,
$r_1=R/8$ (for legs with $a_{j+1}\ge b^*$ they are empty). Theorem~\ref{thm:R}, started at $\Bjs$ (Remark~\ref{rem:B-R}(2) for the
variant), has
\[
16e^{2\mu}+20c_0+16e^{2\mu}\BLs/R+K_T\le32+\tfrac1{1960}+32\BLs/R+K_T=K^*
\]
and $\mu K^*\le v(1+2v)K^*\le0.51$, hence the time is $\le\Bts+\BLs/((1-v)(1-\mu K^*))\le\BTs+2.1\BLs$. For exact capture apply
Lemma~\ref{lem:A41} with any $\varepsilon_N\downarrow0$, $\varepsilon_N\le\varepsilon_0$, $H_N$ the path for $(v,\varepsilon_N)$
(which is $1$-Lipschitz, starts at $0$ and sees all directions; ``seen'' refers to the true candidates by
Lemma~\ref{lem:B-VC}(c)), $T:=\BTs+2.1\BLs$, and arbitrary nested finite sets $D_N$ with dense union.
\end{proof}

\begin{remark}[Status]\label{rem:B-status}
\emph{Refereed.} Lemmas~\ref{lem:B-V}--\ref{lem:B-C}, Theorems~\ref{thm:R} and~\ref{thm:Bulk} and Corollary~\ref{cor:M} (as a
conditional statement) were checked by two independent referees in an earlier version (version~2 of the claims; reports
\texttt{referee-RP6a-1.md} and \texttt{referee-RP6a-2.md}). Neither found a false proved statement; each found one gap (Step~4,
Term~II, case $b<a_{j+1}$; the last leg of Theorem~\ref{thm:Bulk}), and all corrections are incorporated. The corrected passages
(Lemma~\ref{lem:B-C}(0), the scope of Theorem~\ref{thm:Bulk}, the rewritten Step~3 and its constants, the added case in Step~4,
the arithmetic of Steps~2 and~5, Remark~\ref{rem:B-B}(4), the wording of Corollary~\ref{cor:M}) were written after the reports,
some of them in one referee's own words, and were not re-read by the referees. They were re-checked afterwards in an audit of
every difference between the typeset text and the refereed text (report \texttt{audit-B.md}: no difference is false or leaves a
gap; its two corrections to the text of Step~3 are incorporated). Remark~\ref{rem:B-relaxed} was added afterwards and is
covered by no referee report. The parts of this section that are used for Theorems~\ref{thm:Main} and~\ref{thm:Main2} (the lemmas, and the
computations of the proof of Theorem~\ref{thm:Bulk} quoted in the proof of Theorem~\ref{thm:End}) were also audited end to end
as typeset (report \texttt{audit-A.md}; see Remark~\ref{rem:T2-status}).

\emph{Proved and not proved.} Hypothesis \textup{(ST)} of Corollary~\ref{cor:M} is now \textbf{proved} for the start S of
Section~\ref{sec:start}, uniformly in $\varepsilon$ (Theorem~\ref{thm:ST}; that start is covered by the last clause of
Corollary~\ref{cor:M}, and (N2) below explains why construction C's own start does not serve). For Corollary~\ref{cor:M} only \textup{(T)} remains: it is
\textbf{not} proved, so this corollary gives no unconditional range of $v$ (Corollary~\ref{cor:ST}). The unconditional results
of these notes, Theorems~\ref{thm:Main} and~\ref{thm:Main2}, avoid \textup{(T)} altogether. Both are statements about
quantities determined by $(v,\varepsilon)$ through the path.
\textup{(T)} is not a mild condition: once $a_{j+1}>2\pi-\delta$ it bounds the full quantity $\tilde K$ on the closing arc, so the
whole closing phase of the sweep lies inside \textup{(T)}.

\emph{NUMERICAL (evidence only; none of it verifies \textup{(ST)} or \textup{(T)}).}
(N1) For the head-on start, with $v\in\{0.05,0.1,0.2\}$ and $0.005\le\varepsilon\le0.02$, the bounds (IC1) and (IC2) were observed from the
second leg after first contact with $\kappa^2\le0.038$ on arcs up to $32R=0.11$, and $K_{\mathrm{tip}}\le0.27$ from then on. This does not
verify \textup{(ST)}: these values give $\kappa R=6.7\cdot10^{-4}>1/1584$, so (IC3) fails; (IC3) would need
$R\le1/(1584\kappa)\approx3.2\cdot10^{-3}$ and $\varepsilon\le R/64\approx5\cdot10^{-5}$, far below the tested $\varepsilon$; and
the tested $v$ are more than ten orders of magnitude above $10^{-12}$.

(N2) \textup{(ST)} at a \emph{fixed} leg index is \textbf{not} $\varepsilon$-uniform. The first contact leaves a rounded corner of
total turning $\approx2\arcsin v$ at label $0$, of width $\sim\varepsilon$ and curvature $\approx1.7\,v/\varepsilon$, with a tail ahead
of the base that reaches further, in units of $\varepsilon$, as $\varepsilon\to0$. Both referees measured $\sup E/\lambda^3$ (the
value $\kappa^2$ needed in (IC1)) at leg $2$, with independent codes:
\begin{center}
\begin{tabular}{llllll}
\toprule
\multicolumn{6}{l}{Report RP6a-1: $v=0.05$, pairs ahead of the base with $\lambda\le8\varepsilon$}\\
$\varepsilon$ & $4\cdot10^{-3}$ & $10^{-3}$ & $2.5\cdot10^{-4}$ & $6\cdot10^{-5}$ & $1.5\cdot10^{-5}$\\
$\sup E/\lambda^3$ & $0.033$ & $0.027$ & $0.021$ & $0.114$ & $2.50$\\
\midrule
\multicolumn{6}{l}{Report RP6a-2: $v=0.1$, pairs ahead of the base with $\lambda\ge\varepsilon/2$, within arclength $0.11$}\\
$\varepsilon$ & $5\cdot10^{-3}$ & $10^{-3}$ & $2\cdot10^{-4}$ & $4\cdot10^{-5}$ & $10^{-5}$\\
$\sup E/\lambda^3$ & $0.019$ & $0.019$ & $0.40$ & $11.2$ & $190$\\
\bottomrule
\end{tabular}
\end{center}
(background value $0.019$ at $v=0.1$). At later legs the bound is restored: at leg $8$ the first referee
finds $0.012$--$0.034$ for all tested $\varepsilon$, and at leg $5$ the second finds $0.019$--$0.020$ for $\varepsilon\ge4\cdot10^{-5}$. So
$\Bjs$ must grow as $\varepsilon\to0$ (which Corollary~\ref{cor:M} permits), and proving \textup{(ST)} requires a quantitative
description of this corner tail over the first $\Bjs(v,\varepsilon)$ legs, which Theorem~\ref{thm:Bulk} does not cover; the relaxed
start of Remark~\ref{rem:B-relaxed} does not remove the problem. For fixed $v>0$, \textup{(ST)} at $\Bjs=0$ with absolute $\kappa,R$
is false for small $\varepsilon$: (IC1) forces $\kappa\gtrsim0.35\,v/\varepsilon$. The same corner reappears in \textup{(T)} when the
sweep closes at label $2\pi$.

(N3) In-range simulation of construction C by the first referee: $v=10^{-12}$, $\varepsilon=4.5\cdot10^{-5}$, head-on start,
$\Bts=t_2$, labels $[0,0.25]$, about $3200$ legs, with $\kappa=0.2041$ and $R=1/(1584\kappa)=0.00309$ computed from the data at
$\Bts$ (so $\varepsilon=0.93\,R/64$). All hypotheses of Theorem~\ref{thm:Bulk} hold, and (I1), (I2), \eqref{eq:B-CA},
\eqref{eq:B-star}, Step~1(a) and \eqref{eq:B-TL} hold at every leg start and mid-leg: $\max l/\lambda=1.0000000$ against $2$;
$\max E/(\frac12M^2\lambda^{3/2})=0.0156$, i.e.\ a margin of $64$ in (I2); $\min c/(\frac12\min(\lambda,32R))=2.00$; and a margin of $420$
in \eqref{eq:B-TL}. This is a weak test, since at $v=10^{-12}$ the front hardly moves.
\end{remark}
